%%%%%%%%%%%%%%%%%%%%%%%%%%%%%%%%%%%%%%%%%%%%%%%%%%%%%%%%%%%%%%
%%%%%%%% MA-0477 Ecuaciones diferenciales parciales aplicadas %%%%%%%%%
%%%%%%%%%%%%%%%%%%%%%%%%%%%%%%%%%%%%%%%%%%%%%%%%%%%%%%%%%%%%%%
\newpage
\subsection*{MA-0477 Ecuaciones diferenciales parciales aplicadas}
\addcontentsline{toc}{subsection}{MA-0477 Ecuaciones diferenciales parciales aplicadas}

\hypertarget{MA0477}{}
\begin{refsection}
\begin{table}[H]
\centering{
\begin{tabular}{|ll|}
\hline
\textbf{Año:} IV  & \textbf{Requisitos:} MA-0515 Ecuaciones Diferenciales; MA-0552 Principios de Análisis II \\ 
\textbf{Ciclo:} Optativa  & \textbf{Correquisitos:} Ninguno \\ 
\textbf{Tipo de curso:} Teórico & \textbf{Horas presenciales por semana:} 5 \\ 
\textbf{Créditos:} 4 & \textbf{Horas de trabajo independiente por semana:} 7 \\
\textbf{Virtualidad:} P  &  \\ \hline
\end{tabular}
}
\end{table}

\subsubsection*{Descripción del curso}

Este curso optativo introduce las ecuaciones diferenciales parciales desde un enfoque clásico y aplicado, centrado en métodos explícitos, principios del máximo, funciones de Green y representaciones integrales. Está orientado a la persona estudiante que requiere un dominio sólido de las EDP lineales más usuales en modelación y en la continuación hacia cursos de análisis numérico o estadística aplicada, sin desarrollar por completo la teoría abstracta de espacios de Sobolev.

El programa recorre EDP de primer orden, problemas elípticos (Laplace y Poisson), la ecuación del calor y la de onda, y cierra con una introducción al cálculo de variaciones como puente hacia soluciones débiles. 

\subsubsection*{Objetivos}

\textbf{General:} Aplicar métodos clásicos del análisis de EDP para resolver y analizar ecuaciones lineales que surgen en modelos matemáticos.

\textbf{Específicos:} Al finalizar el curso se espera que la persona estudiante sea capaz de:

\begin{enumerate}
\item Resolver EDP de primer orden mediante características y reconocer formulaciones de leyes de conservación.
\item Utilizar funciones de Green, fórmulas del valor medio y principios del máximo en problemas de Laplace y Poisson.
\item Obtener y aplicar soluciones fundamentales del calor, el principio de Duhamel y principios del máximo o del valor medio asociados.
\item Aplicar la fórmula de d'Alembert, medias esféricas y métodos de energía en la ecuación de onda.
\item Formular problemas variacionales sencillos y relacionarlos con soluciones débiles sin desarrollar la teoría completa de Sobolev.
\item Comunicar resultados y estrategias de resolución con claridad matemática.
\end{enumerate}

\subsubsection*{Contenidos}

\begin{enumerate}

\item \textbf{EDP de primer orden (3 semanas):}
\begin{enumerate}
    \item Ecuación de transporte y método de características.
    \item Introducción a leyes de conservación.
\end{enumerate}

\item \textbf{Ecuaciones de Laplace y Poisson (4 semanas):}
\begin{enumerate}
    \item Funciones de Green y fórmulas del valor medio.
    \item Principios del máximo y fórmula de Poisson.
\end{enumerate}

\item \textbf{Ecuación del calor (3 semanas):}
\begin{enumerate}
    \item Solución fundamental y principio de Duhamel.
    \item Principios del valor medio y del máximo.
\end{enumerate}

\item \textbf{Ecuación de onda (3 semanas):}
\begin{enumerate}
    \item Fórmula de d'Alembert y medias esféricas (fórmulas de Kirchhoff y de Poisson).
    \item Métodos de energía.
\end{enumerate}

\item \textbf{Introducción al cálculo de variaciones (2 semanas):}
\begin{enumerate}
    \item Minimizadores de funcionales de energía y soluciones débiles.
\end{enumerate}

\end{enumerate}

\subsubsection*{Metodología}

\noindent \textbf{Lineamientos metodológicos:} La persona docente a cargo de este curso debe respetar las siguientes pautas:

\begin{enumerate}
    \item Combinar \textbf{clases magistrales} con \textbf{sesiones de ejercicios} orientadas a cálculo de soluciones explícitas y aplicación de principios cualitativos.
    \item Enfatizar ejemplos y modelos que motiven cada familia de EDP.
    \item Asignar tareas regulares que integren teoría y manipulación de fórmulas clásicas.
\end{enumerate}

\textbf{Sugerencias metodológicas:} Adicionalmente, se tienen las siguientes recomendaciones para la persona docente a cargo de este curso:

\begin{enumerate}
    \item Señalar la articulación natural con MA-0756 (teoría abstracta) y MA-0920 (aproximación numérica).
    \item Reservar la última semana para conectar el cálculo de variaciones con la noción de solución débil vista en cursos posteriores.
    \item Fomentar la comunicación oral breve de estrategias de resolución en la pizarra.
    \item El tema de cálculo variacional se puede hacer sin introducir la maquinaria completa de espacios de Sobolev
\end{enumerate}

\nocite{Eva97}
\nocite{john82}
\nocite{strauss07}
\nocite{weinberger95}

\printcursosbibliography

\end{refsection}
